\chapter{Технологическая часть}
В данной части будут описаны требования к реализации программного обеспечения, выбору языков программирования и тестированию.


\section{Требование к программному обеспечению}
Программное обеспечение должно удовлетворять следующим критериям:
\begin{itemize}
	\itemпрограмма должна обеспечить возможность ввести последовательность для последующего нахождения максимума её элементов. Для успешного ввода пользователь должен иметь возможность задать длину последовательности и  поэлементно ввести содержимое матрицы;
	\itemдля замеров времени выполнения, ПО должно обеспечить возможность ввода начального и конечного размера матриц, а также ввести шаг замеров;
	\itemпрограмма должна иметь возможность нахождения максимума двумя алгоритмами: итерационным и рекурсивным
\end{itemize}

\section{Средства реализации}
В качестве средства реализации был выбран язык C++[5], из-за возможности управления памятью и наличия контейнерных классов, встроенных в стандартную библиотеку . Разработка проводилась в среде программирования VS Code. Для сборки использована утилита CMake. Для визуализации данных, полученных в ходе исследования использовался язык Python[6] за счёт обилия библиотек для анализа данных. Мною были использованы pandas[7] и matplotlib[8]. Для запуска тестов и последующего сохранения результатов также использовался язык Python.

\section{Реализация алгоритмов}
При выполнении лабораторной работы были реализованы два варианта алгоритма поиска максимума. Реализация итерационного алгоритма указан на листинге \ref{list:iter}. Реализация рекурсивного алгоритма указан на листинге \ref{list:rec}
\\
\begin{lstlisting}[language=C++, caption={Итерационный алгоритм поиска максимума}, captionpos=b, label=list:iter]
void Sequence::default_max(int &max) {
	max = INT_MIN;
	for (size_type i = 0; i < _len; i++) {
		if (_container[i] > max)
		max = _container[i];}}
\end{lstlisting}

\clearpage
\begin{lstlisting}[language=C++, caption={Рекурсивный алгоритм поиска максимума}, captionpos=b,  label=list:rec]
void Sequence::recursion_max(int &max, size_type index) {
	if (index == _len)
		return;
	if (_container[index] > max)
		max = _container[index];
	recursion_max(max, index += 1);}
\end{lstlisting}

\section{Тестирование ПО}
Для тестирование реализованы ПО были выделены следующие тестовые случаи для каждого алгоритма:
\begin{table}[h]
	\centering
	\caption{Тестовые случаи для функции рекурсивного поиска максимума}
	\begin{tabular}{|c|c|c|c|}
		\hline
		\textbf{Входная последовательность} & \textbf{Ожидаемый максимум} &  \textbf{Описание теста} \\
		\hline
		\(\{-1,\ -3,\ 0\}\) & \(0\) & \text{Максимум ноль} \\
		\hline
		\(\{-1,\ -3,\ -6\}\) & \(-1\)  & \text{Максимум отрицательный}  \\
		\hline
		\(\{100,\ 23,\ 6\}\) & \(100\) & \text{Максимум положительный}  \\
		\hline
		\(\{100\}\)  & \(100\) & \text{Последовательность из 1 элемента} \\
		\hline
		\(\{100,\ 6\}\)  & \(100\) & \text{Последовательность из 2 элементов} \\
		\hline
		\(\{100,\ 50,\ 20,\ 40\}\) & \(100\) & \text{Максимум первый} \\
		\hline
		\(\{100,\ 50,\ 33,\ 200\}\)& \(200\) & \text{Максимум последний} \\
		\hline
		\(\{40,\ 20,\ 10,\ 0,\ 100,\ 50\}\) & \(100\) & \text{Максимум в середине}  \\
		\hline
	\end{tabular}
\end{table}

Все тесты прошли успешно.

\section*{Вывод}
На основе конструкторской части были реализованы и протестированы алгоритмы нахождения максимума матрицы.

\clearpage
